% ============================================================
%  CAPÍTULO 6 – CONCLUSIONES Y TRABAJO FUTURO
% ============================================================

\subsection{Conclusiones}
\label{subsec:conclusiones}

Este trabajo ha desarrollado CAIQ (\textit{Career Alignment and Insight Qualifier}),
un sistema de recomendación profesional orientado a personas en transición hacia el
mercado laboral de datos. El sistema toma como entrada el currículum vitae del candidato,
cuantifica su brecha de competencias respecto a un rol objetivo y genera recomendaciones
ordenadas de formación y empleo. El resultado cubre el ciclo completo de un proyecto de
datos: desde la adquisición y preprocesamiento de fuentes heterogéneas hasta el despliegue
de una aplicación web de acceso público con actualización automática semanal del corpus
subyacente.

\textbf{Cobertura de los objetivos.} Los cinco objetivos específicos planteados quedan
cubiertos con evidencia cuantitativa. El sistema dispone de tres \textit{datasets}
estructurados y enriquecidos (37.829 ofertas de empleo, 89.332 cursos y 3.503 másteres)
mantenidos mediante un \textit{pipeline} incremental con deduplicación, caducidad
temporal y verificación activa de URLs. El motor de recomendación ha sido comparado en
cuatro variantes mediante ablación sobre $n$\,=\,482 candidatos y validado tanto
cuantitativamente como sobre casos de uso representativos. La aplicación está desplegada
y operativa en Hugging Face Spaces. El detalle de los resultados y su interpretación
crítica se ha desarrollado en los capítulos anteriores.

\textbf{Contribuciones principales.} El trabajo aporta tres contribuciones con valor
propio más allá de los objetivos. En primer lugar, una \textbf{taxonomía controlada de
71 competencias del sector de datos} con 477 alias léxicos, construida sobre análisis
empírico de ofertas reales, reutilizable como recurso independiente para otros sistemas
de análisis de mercado laboral. En segundo lugar, un \textbf{motor de recomendación
híbrido sin dependencia de historial de interacciones}: al no requerir datos previos de
uso, el sistema ofrece recomendaciones desde la primera consulta, eliminando el problema
de arranque en frío que limita los sistemas de filtrado colaborativo. En tercer lugar,
un \textbf{\textit{pipeline} de adquisición y mantenimiento incremental} que combina
scraping multi-fuente, verificación activa de URLs con caché persistente y caducidad
configurable, reproducible en producción con infraestructura mínima.

\subsection{Trabajo Futuro}
\label{subsec:trabajo_futuro}

Los resultados obtenidos y las limitaciones identificadas permiten trazar líneas de
trabajo concretas que ampliarían la capacidad y el alcance del sistema.

\textbf{Mejora de la cobertura del catálogo de cursos.} La limitación más directa sobre
la calidad de las recomendaciones formativas es la baja cobertura de competencias en los
cursos (15,7\,\%), consecuencia de disponer solo del título como fuente textual. La
prioridad inmediata es enriquecer el catálogo con descripciones de contenido y
\textit{syllabi} obtenidos de las APIs o páginas de cada plataforma. Esta mejora no
requiere cambios arquitectónicos y tendría impacto directo en la precisión del ranking
formativo.

\textbf{Extracción de competencias con modelos preentrenados.} La taxonomía léxica
actual ofrece cobertura casi total pero precisión moderada. Una vez disponible
infraestructura de inferencia con GPU, sustituir o complementar el extractor actual con
modelos especializados en vocabulario de empleo --- como ESCOXLM-R (Zhang et al., 2023),
preentrenado sobre la taxonomía ESCO --- permitiría reducir los falsos positivos sin
sacrificar cobertura.

\textbf{Incorporación de \textit{feedback} implícito del usuario.} El sistema genera
recomendaciones sin memoria de sesiones anteriores. Registrar las interacciones del
usuario (cursos consultados, ofertas guardadas) como señal de relevancia implícita
permitiría recalibrar los pesos del motor de forma continua, añadiendo una capa de
personalización basada en comportamiento real sin necesidad de valoraciones explícitas.

\textbf{Expansión y mantenimiento dinámico de la taxonomía.} La taxonomía de 71
competencias cubre el mercado actual del sector de datos, pero el vocabulario técnico
evoluciona con rapidez. Definir un proceso de revisión periódica anclado en el análisis
automático de términos emergentes en las ofertas scrapeadas --- identificando denominaciones
frecuentes no mapeadas --- permitiría mantener la taxonomía alineada con la demanda real
sin intervención manual intensiva.

\textbf{Evaluación con usuarios reales.} La validación del sistema se ha realizado sobre
CVs reales pero sin interacción directa de candidatos con la aplicación. Un estudio con
profesionales en transición que utilizasen CAIQ durante un período definido permitiría
medir la utilidad percibida, la relevancia subjetiva de las recomendaciones y el impacto
sobre las decisiones de formación, complementando la evidencia cuantitativa con la
perspectiva del usuario final.

\textbf{Adaptación a entornos organizacionales.} Las aplicaciones identificadas en la
discusión --- servicios de orientación universitaria y departamentos de recursos humanos ---
requieren extender el sistema con gestión de usuarios, soporte para análisis colectivos
(\textit{team gap analysis}) y exportación de informes estructurados. El núcleo analítico
actual podría actuar como motor de estas extensiones sin necesidad de rediseño.

En conjunto, CAIQ demuestra que es posible construir un sistema de orientación profesional basado en datos reales del mercado ---funcional, validado y desplegado--- con infraestructura accesible y sin dependencia de modelos de lenguaje de gran escala. El profesional en transición dispone así de una referencia objetiva sobre su punto de partida y de los pasos concretos que puede dar para cerrar la brecha que le separa del rol al que aspira.
